\section{Introduction}
\label{sec:intro}

The rapid expansion of artificial intelligence (AI) is turning data centers into one of the fastest-growing
electricity loads. The International Energy Agency estimates that data centers consumed about 415~TWh in 2024, roughly
1.5\% of global electricity use, and that their demand may reach about 945~TWh by 2030~\cite{iea2025}. In the United
States their share already reached 4.4\% in 2023 and may rise to 325--580~TWh by 2028~\cite{shehabi2024}, while China
requires new large data centers to achieve a power usage effectiveness (PUE) of 1.25 or lower~\cite{ndrc2024dc}.
Because new AI campuses are increasingly limited by grid interconnection rather than by chips, operational flexibility
is becoming a condition of connection: curtailing large flexible loads for 0.25--1\% of the hours could accommodate
76--126~GW of additional load in 22 balancing authorities of the United States~\cite{norris2025}, the flexibility of AI loads can reduce grid
investment and operating costs by 3--21\%~\cite{chen2026defer}, and the Texas Senate Bill~6 links the interconnection of
large loads to curtailment obligations during grid emergencies~\cite{texas2025sb6}. Power-flexible AI data centers are
therefore regarded as a new class of grid-responsive resources~\cite{williams2026}, and a field demonstration in Phoenix
reduced the power of a 256-GPU cluster by 25\% for three hours~\cite{colangelo2025}. In China, demand-side regulations
require a demand response (DR) capability of 3--5\% of the maximum load~\cite{ndrc2023dsm}. Under the Jiangsu DR rules,
an event is valid only if the mean load falls by at least 80\% of the confirmed reduction, payments are capped at
120\% of it, agreed responses are limited to two calls and 2~h per day, and fast responses with short notice are paid up to 1.6 times
the base rate~\cite{jiangsu2024dr}. One evening peak can therefore trigger several events that are separated by gaps of
an hour or less, and an AI data center that responds to the first event must be ready for the next one shortly
afterwards.

In the area of data center DR and AI workload flexibility, early studies showed that data centers can avoid coincident
peaks by shifting deferrable workloads and using local generation~\cite{liu2013}, discussed the opportunities and
challenges of data center DR~\cite{wierman2014}, and reviewed how Internet data centers participate in DR
programs~\cite{chen2020review}. Carbon-aware computing shifts flexible workloads towards low-carbon hours at fleet
scale~\cite{radovanovic2023}. For AI workloads, simulation studies estimate a power flexibility of 18--55\% for AI data
centers~\cite{acun2026}, power capping and frequency scaling provide additional levers for large language model (LLM)
workloads~\cite{patel2024}, and quantization can turn LLM inference into a DR resource~\cite{du2026}. Agreements that
curtail machine-learning workloads during grid stress have also been signed in practice~\cite{google2025}. However, a
study of 185 days of production traces from 155,000 GPUs found that a scalar flexibility ratio overstates the dependable
load relief by 17--47\%~\cite{li2026scalar}, and synchronous training causes large power swings~\cite{choukse2025}.
These studies characterise how much an AI data center can reduce in a single event, but they do not consider how the
recovery needs left by one event constrain the next event.

In the area of data center integrated energy systems (IES), multi-level DR dispatch of data centers with district
autonomy~\cite{han2024}, exergy-based planning and operation of computing-power IES~\cite{lin2025} and the operation
of data centers with combined cooling, heating and power (CCHP) and spatio-temporally adjustable loads~\cite{fan2026}
have been studied. Zhang et al.~\cite{zhang2026bilevel} proposed a bi-level distributionally robust planning model for
data center clusters in which time-shiftable and space-transferable loads act as virtual energy storage, and the reuse of
data center waste heat in district heating has been reviewed~\cite{huang2020,wahlroos2018}. In these models, deferred
workloads are bounded by constant virtual-queue limits and recovered by their own rules, while storage is recharged by
the energy system, so the recovery of computing backlog and that of storage are treated as separate subproblems. As
shown in \secref{sec:computing}, a constant queue limit cannot guarantee that every task meets its deadline when
arrivals are low and is unnecessarily conservative when arrivals are high.

In the area of recovery and coordination, Evans et al.~\cite{evans2022} proposed a flexibility framework with recovery
guarantees for aggregated energy storage, in which discharge and recovery are coupled and the recovery capability is
summarised by a single virtual battery; virtual-battery models also describe the aggregate flexibility of
thermostatically controlled loads~\cite{hao2015}. Load payback after DR events is a common side effect of
DR~\cite{wrinch2012}, and the choice of the customer baseline affects the settled reduction~\cite{chao2011}.
Uncertainty in DR scheduling is usually handled by scenario-based stochastic model predictive control
(MPC)~\cite{schildbach2014} or by distributionally robust optimization (DRO), which in two-stage problems is typically
solved by duality-based reformulation and column-and-constraint generation~\cite{zhang2026bilevel,zeng2013,esfahani2018}.
Coordination between the IT scheduler and the facility energy management system is commonly realised by price signals,
whereas iterative schemes such as the alternating direction method of multipliers (ADMM)~\cite{boyd2011} exchange both
plans and prices. However, existing approaches consider either one homogeneous storage state or a single event, and
none addresses two heterogeneous recovery states, a deadline-bounded and nearly lossless computing backlog and a lossy
physical storage, that compete for one connection-limited recovery channel across consecutive events.

Despite substantial progress in data center DR, IES operation and storage recovery, several research gaps remain.
First, existing studies optimise a single event or treat events independently, so the state coupling between
consecutive events, \ie\ the computing backlog and the storage deficit that must be recovered through the same
connection headroom, is not modelled, and the resulting loss of deliverable capability is unknown. Second, computing
backlog and storage are recovered by separate rules or coordinated by static prices, whereas their joint value, the
conditions under which this value vanishes and the price that would coordinate them have not been characterised.
Third, commitments are assessed event by event and against weak benchmarks, without a description of the set of
deliverable commitment vectors, out-of-sample validation over many operating days, or a job-level verification of the
fluid workload model. Consequently, a scheduling framework that jointly recovers computing and storage across events,
from the day-ahead commitment to intraday operation, is still lacking.

To address these gaps, this paper proposes a coordinated scheduling method for consecutive DR in AI data center IES
(AIDC-IES) considering cross-event computing--storage recovery. The main contributions are summarised as follows.
\begin{enumerate}[label=(\arabic*),leftmargin=1.9em,itemsep=1.5pt]
\item A multi-state recovery model of computing and storage is established. Tiered service-level agreements (SLAs) are
represented by lag-bounded backlog constraints, which are proved to be exactly equivalent to deadline-feasible
earliest-deadline-first (EDF) execution (Lemma~\ref{lem:edf}), so that deferrable computing becomes a virtual storage
whose capacity is set by deadlines. Together with the battery, the chilled-water storage, the gas engine with
absorption cooling and the connection limit, the model captures how catch-up computing, recharging and the next
event's capability interact; data hall thermal inertia, waste-heat recovery, GPU power capping, pause overheads and a
rigid-job recovery bound are included as options.
\item Theoretical properties of cross-event recovery are derived: an upper bound on the deliverable capability of the
next event (Proposition~\ref{prop:bound}), a closed-form value of joint recovery and of the minimum recovery interval
(Theorem~\ref{thm:closed}), a single recovery-channel price with equal marginal values
(Proposition~\ref{prop:price}), an exact mixed-integer linear reformulation of the distributionally robust commitment
(Proposition~\ref{prop:dro}), the convex capability region of commitment vectors with a virtual-battery description
and a vertex certificate (Proposition~\ref{prop:region}), and a distributed plan--price iteration that converges to
joint recovery and recovers the channel price (Proposition~\ref{prop:admm}).
\item A day-ahead coupled commitment and intraday rolling joint-recovery framework is developed
(Algorithm~1, \tabref{tab:algorithm}). It embeds the Jiangsu settlement rules, handles uncertain event calls by
scenario trees, and admits a distributed implementation in which the IT scheduler and the facility controller keep their
own models and exchange only power plans and prices.
\item A comprehensive evaluation against strong benchmarks is conducted on a 16,000-GPU case built from public
traces. The benchmarks include independent recovery with a back-tested split, independent recovery coordinated by the
channel price of the joint plan and event-by-event joint recovery. The evaluation covers 52 out-of-sample operating
days with paired tests, risk-averse commitments, the capability region, job-level replay of all schedules and
extensions to green-power direct supply and winter waste-heat recovery, and it quantifies the reliably deliverable
capacity, its equivalent in battery and connection capacity, and the conditions under which joint recovery loses its
value.
\end{enumerate}
